\chapter{Backtesting Done Right}
\label{ch:backtesting}
\chaptermark{Backtesting done right}

A \textbf{backtest} is a simulation of what a trading strategy would
have made on historical data. Done wrong, it is a lie: a strategy
that ``made 20\% per year in backtest'' can lose money live, for
avoidable reasons. \Cref{ch:mean_reversion_ou,ch:pairs_cointegration}
already met two (the look-ahead from fitting an OU model on the
full dataset, and the over-refit hedge ratio), and in both the
cure was walk-forward. This chapter treats backtest discipline as a
first-class subject and catalogues the \emph{seven standard ways
backtests lie}, with numerical examples and defences. The centrepiece
is Lop\'ez de Prado's \textbf{Combinatorial Purged Cross-Validation}
(CPCV), the minimum cross-validation discipline for financial machine
learning.

\Cref{ch:kelly_risk} scored track records with Sharpe, Sortino,
Calmar, information ratio, and drawdown. Every one of those numbers
\emph{only means something if the backtest was honest}. A Sharpe of
$1.8$ from $1000$ parameter sweeps is worth about $0.6$ after
multiple-testing deflation, as derived by hand in
\cref{sec:bt_deflated_sharpe}. Hence the \textbf{Deflated Sharpe
Ratio} \citeAFML{} as the principled correction.

Read this chapter as an anti-handbook: most of the value is never
making these mistakes, the rest is knowing which framework
(\texttt{Backtesting.py}, \texttt{vectorbt}, \texttt{NautilusTrader},
\texttt{QuantConnect}) to reach for.

\begin{backgroundbox}[Prerequisites]
From earlier in the book:
\begin{itemize}
  \item \Cref{ch:mean_reversion_ou}: walk-forward OU fitting, and
  the pitfall of computing historical $z$-scores against an
  $\hat\mu$ estimated on the full series. That chapter's
  pitfall on look-ahead OU fitting is the canonical example of
  Lie~1 below.
  \item \Cref{ch:pairs_cointegration}: the ``over-refit
  $\hat\beta$'' pitfall (refitting the hedge ratio every bar
  on windows that include the current observation) is the
  canonical example of Lies~1 and~3 acting together.
  \item \Cref{ch:kelly_risk}: the Sharpe ratio, the drawdown
  recovery identity, and the framing of performance metrics as
  ``scoring after sizing''. This chapter's business is to ensure
  those metrics are not meaningless.
  \item \Cref{ch:market_impact} and~\cref{ch:almgren_chriss}:
  the square-root impact law and the A--C execution cost model.
  Lie~5 (transaction-cost neglect) is resolved using numbers
  from those chapters.
\end{itemize}
From outside the book:
\begin{itemize}
  \item Elementary statistics: $t$-test, $p$-value, standard
  error, and the idea of a hypothesis test under a null.
  \item Machine-learning cross-validation: $k$-fold CV as it is
  taught in any intro ML course. This chapter explains why the
  textbook version is wrong for financial time series, and
  what to replace it with.
  \item The \textbf{expected maximum of $N$ i.i.d.\ standard
  normals} is approximately $\sqrt{2 \ln N}$; this heuristic
  controls the magnitude of data-snooping bias in
  \cref{sec:bt_deflated_sharpe}.
\end{itemize}
\end{backgroundbox}

% =====================================================================
\section{The seven lies a backtest can tell you}
\label{sec:bt_seven_lies}

Every subsequent section (walk-forward, CPCV, deflated Sharpe,
framework choice) defends against one or more of these seven
failure modes. The list is adapted from \citet{lopezdeprado2018afml}
(Ch.~11--12) and \citet{chan2013algorithmic} (Ch.~3).

\subsection{Lie 1: Look-ahead bias}
\label{subsec:lie_lookahead}

A backtest commits \textbf{look-ahead bias} when the simulated
strategy, at time $t$, uses information unavailable at $t$ in real
trading. The cartoon version: the backtest buys today because it
already knows tomorrow's closing price went up; live, tomorrow's
close does not yet exist. Archetype: the OU pitfall of \cref{ch:mean_reversion_ou},
where $\hat\mu$ and $\hat\theta$ are fitted on the \emph{full} series
and used to compute every historical $z$-score. Each entry is timed
against a mean estimated using post-trade data. Live, $\hat\mu$ is
fitted on a rolling window ending before $t$ and differs meaningfully
from the in-sample value, so the thresholds hit in backtest vanish.

Most sources do not look like look-ahead:
\begin{itemize}
  \item \textbf{Close-to-close returns used intraday.} A signal
  that reacts to $r_t = \log(C_t / C_{t-1})$ at time
  $t + \epsilon$ has used the closing print from time $t$, which
  was not available until the closing auction completed.
  \item \textbf{Corporate-action adjustments.} Dividend- and
  split-adjusted historical prices reflect \emph{future}
  corporate actions: if ACME paid a \$$2$ dividend in 2024 and
  you pulled adjusted prices in 2026, the 2022 prices are
  shifted down by the dividend. A 2022-live strategy had no
  access to that adjustment.
  \item \textbf{Fundamental data with revision lag.} Earnings,
  GDP, index reconstitution: the value you see in your data
  vendor for date $t$ is the \emph{latest revision}, not the
  value the market saw on $t$.
  \item \textbf{Point-in-time benchmarks.} The S\&P~500
  constituents as of a 2015 date are not the current
  constituents; the historical series you downloaded may be.
  This is the gateway to Lie~2.
\end{itemize}

\paragraph{Concrete magnitude.} OU pairs trades with $\hat\mu$ on the
full sample report Sharpes of $2.0$--$3.0$ on canonical equity pairs;
switching to strict walk-forward ($\hat\mu$ on a rolling window ending
at $t - 1$) typically cuts the Sharpe to $0.5$--$1.0$ on the same
series \citeChan{}. The difference is pure look-ahead.

\subsection{Lie 2: Survivorship bias}
\label{subsec:lie_survivorship}

A backtest commits \textbf{survivorship bias} when run on a universe
defined by \emph{today's} membership. Example: a long-only
equal-weighted S\&P~500 backtest $2000$--$2020$ using current
constituents. The current list contains firms that \emph{survived};
bankruptcies, delistings, and distressed acquisitions between 2000
and 2020 are missing. Enron, Lehman, WorldCom, Washington Mutual,
pre-bankruptcy GM, Bear Stearns: none appear, because none are in
today's constituent list.

\paragraph{Magnitude.} \citet{lopezdeprado2018afml} and
\citet{chan2013algorithmic} both report that survivorship inflates
long-only S\&P~500 backtests by $0.2$--$0.5$ points of annualised
return per year on a 20-year window: a Sharpe inflation of
$0.3$--$0.5$ at $15\%$ vol. Tighter universes are worse: long/short
on micro-cap NASDAQ using a \emph{current} universe can gain
$2$--$5$ Sharpe points from survivorship, because the delisted
micro-caps were catastrophic losers.

The only fix is a \textbf{point-in-time universe}: tradable
instruments \emph{as known on date $t$}, including now-delisted
names, with ticker changes and delisting events. CRSP, Bloomberg,
and QuantConnect's Lean dataset ship PIT universes; free sources
typically do not. Assume a free dataset is survivorship-biased
unless documented otherwise.

\subsection{Lie 3: Overfitting}
\label{subsec:lie_overfit}

A backtest commits \textbf{overfitting} (\emph{in-sample tuning})
when the same data is used to both \emph{select} parameters and
\emph{evaluate} performance. Archetype: a grid search over $1000$
(lookback, $z$-threshold, stop-loss) combinations on the same 5-year
returns, reporting the best Sharpe. Under a zero-edge null the
winner's Sharpe is about $\sqrt{2 \ln N}$ standard errors above zero:
for $N = 1000$, $T = 1260$, this is
$\sqrt{2 \ln 1000} / \sqrt{T / 252} \approx 1.66$ annualised.
Random strategies grid-searched $1000$ times routinely produce
Sharpes near $1.66$ with no alpha.

Overfitting is distinct from Lie~4 (data snooping): overfitting tunes
\emph{one} strategy's parameters on the evaluation data; snooping
tries many \emph{distinct} ideas and reports only winners. Both
inflate Sharpe by multiple testing; the defence is CPCV
(\cref{sec:bt_cpcv}) plus Deflated Sharpe
(\cref{sec:bt_deflated_sharpe}).

\subsection{Lie 4: Data snooping / selection bias}
\label{subsec:lie_snooping}

A backtest commits \textbf{data snooping bias} (selection bias) when
the researcher evaluates many distinct strategy ideas on the same
dataset and reports only the winners. The winner's Sharpe is biased
upward by the same multiple-testing inflation as Lie~3.

The dangerous case is not counting trials. Six months of iteration
typically touches $N \in [50, 500]$ strategies; uncounted, the
Sharpe in the final writeup is silently inflated and the strategy
goes live on noise. A research log (every idea, parameter set, and
result) is the only mechanical defence;
\cref{sec:bt_deflated_sharpe} turns it into a formula.

\subsection{Lie 5: Transaction-cost neglect}
\label{subsec:lie_txcost}

A backtest commits \textbf{transaction-cost neglect} when it assumes
free execution, fills at the mid, or no market-impact penalty. For a
fast strategy this is the difference between ``gold mine'' and
``money sink''.

\paragraph{Worked magnitude.} Take a daily-frequency mean-reversion
strategy with gross expected per-day return $\mu = 5$~bp and
per-day return standard deviation $\volat = 50$~bp. Its gross
annualised Sharpe is
\[
\text{SR}_{\text{gross}}
  \;=\; \frac{\mu}{\volat}\sqrt{252}
  \;=\; \frac{0.0005}{0.0050}\sqrt{252}
  \;=\; 1.587.
\]
Net of a $2$~bp round-trip cost per trade (half-spread plus small
impact), the mean collapses to $\mu_{\text{net}} = 3$~bp and
$\text{SR}_{\text{net}} = (0.0003/0.0050)\sqrt{252} = 0.952$. Net of
$5$~bp, $\mu_{\text{net}} = 0$ and $\text{SR}_{\text{net}} = 0$.
Edge at $2$~bp, break-even at $5$~bp. A costless backtest reports
Sharpe~$1.59$ in all three cases.

On a sell-side desk, transaction cost combines the square-root
impact law of \cref{ch:market_impact}
($I(Q) \approx Y\volat\sqrt{Q/V}$, $Y \approx 0.1$ for cash equities)
with the half-spread from \cref{ch:spreads_roll} and explicit
fees/rebates. A Prosperity backtest that subtracts a flat tick is a
reasonable first cut; zero is a lie.

\subsection{Lie 6: Slippage}
\label{subsec:lie_slippage}

A backtest commits \textbf{slippage} misspecification when it
assumes execution at the backtest-time price but live execution
happens later, at a worse price, or walks the book. Transaction cost
is the expected cost of crossing the spread and consuming visible
liquidity; slippage is the \emph{unanticipated} wedge between intent
and fill.

Common backtest slippage lies:
\begin{itemize}
  \item \textbf{``Fill at the signal-bar close.''} The decision
  arrives after bar $t$ closes, so the earliest fill is the open
  of $t + 1$, with the full overnight gap in between.
  \item \textbf{``Fill at the mid.''} Market orders fill at the
  best opposing quote. On a symmetric book this is half the
  spread (already in Lie~5); on an asymmetric book, more.
  \item \textbf{``Fill my full size at one level.''} A
  $10{,}000$-lot buy against top-of-book size $100$ walks the
  book; filling the whole order at the top ask underestimates
  cost.
  \item \textbf{``No queue at the top.''} A passive bid does not
  fill just because the market trades there; its price-time
  priority position matters. Assuming instant passive fills is
  over-optimistic by a factor proportional to queue length.
\end{itemize}
The fix is a slippage model: flat $k$ ticks, a linear-in-volume
model calibrated from bid/ask imbalance, or a full order-book
simulation like NautilusTrader or Lean provide.

\subsection{Lie 7: Regime dependence}
\label{subsec:lie_regime}

A backtest commits \textbf{regime dependence} fraud when the period
is dominated by a regime that does not hold live. A trend-follower
profitable $2009$--$2019$ (central-bank-supported trends) dies if
extended back to the $1987$ crash or forward into the $2022$
tightening. A vol seller looks wonderful $2013$--$2019$ and
catastrophic on COVID-March 2020 or ``Volmageddon'' 2018.

Regime dependence is not ``the backtest was wrong'' but ``the sample
was unrepresentative''. There is no statistical cure: statistics
assume the future resembles the past. Partial cures:
\begin{itemize}
  \item Extend the window as far back as data supports.
  $2002$--$2022$ spans four distinct regimes (dot-com aftermath,
  GFC, QE, post-COVID tightening); $2015$--$2020$ covers one
  regime plus a crisis spike.
  \item \textbf{Stratify by regime} using a VIX- or yield-curve-
  based label, reporting Sharpe, return, drawdown per regime. A
  strategy with Sharpe $1.8$ low-vol and $-2.0$ high-vol is two
  strategies that should be switched on a signal.
  \item Reserve a \textbf{final hold-out}, the last $1$--$2$
  years, \emph{never} looked at during research. Evaluate on
  it exactly once, at the end.
\end{itemize}

\begin{pitfallbox}[The seven deadly backtest sins]
\begin{enumerate}
  \item \textbf{Look-ahead}: using data from after decision
  time.
  \item \textbf{Survivorship}: running a universe defined by
  today's membership.
  \item \textbf{Overfitting}: tuning parameters on the same
  data used for evaluation.
  \item \textbf{Data snooping}: trying many strategies on the
  same data and reporting only the winners.
  \item \textbf{Transaction-cost neglect}: omitting spread,
  impact, financing.
  \item \textbf{Slippage}: assuming the fill price equals the
  signal-time price.
  \item \textbf{Regime dependence}: relying on a sample
  period whose conditions do not persist.
\end{enumerate}
Each can turn a break-even strategy into a ``backtested
Sharpe~$2.0$'' that evaporates day one live.
\end{pitfallbox}

% =====================================================================
\section{Walk-forward validation}
\label{sec:bt_walkforward}

The first-order defence against Lies~1 and~3 is \textbf{walk-forward
validation}, met in \cref{ch:mean_reversion_ou} for honest OU
fitting. Here it is a general protocol.

\paragraph{The procedure.} Pick a fitting window $W$ (e.g.\ $1$
year) and a test step $R$ (e.g.\ $1$ month). Run the backtest as
a loop:
\begin{enumerate}
  \item Fit all strategy parameters on observations in
  $[t - W,\; t)$. These are the \emph{training} observations.
  \item Freeze the parameters. Run the strategy on observations
  in $[t,\; t + R)$. These are the \emph{test} observations.
  Record the realised P\&L.
  \item Slide the window: $t \leftarrow t + R$. Go to step~1.
\end{enumerate}
Sharpe, return, and drawdown come from the \emph{concatenation of
test-period P\&Ls}. No test observation was used in fitting the
parameters that generated its signal, so the result is
out-of-sample by construction.

\begin{keyfactbox}[Walk-forward procedure]
Fix window $W$ and step $R$. For each time $t$:
\[
\underbrace{[\,t - W,\, t\,)}_{\text{fit parameters}}
  \quad\Rightarrow\quad
\underbrace{[\,t,\, t + R\,)}_{\text{trade and record P\&L}}.
\]
Slide $t \to t + R$ and repeat. Concatenate the test-period
P\&Ls. Every observation the strategy trades was \emph{unseen}
when its parameters were fit.
\end{keyfactbox}

\paragraph{Why it beats a single split.} Walk-forward simulates the
real passage of time: repeatedly fit on the past, trade the next
unseen period, then advance, which is exactly what a live strategy does.
A random train/test shuffle (the ML default) would scatter future
points into the training set, leaking tomorrow's prices into
yesterday's fit; a fixed split
(``$2000$--$2015$ train, $2016$--$2022$ test'') has two problems.
First, the test period lives in a single regime. Second, the test
Sharpe has high variance: $6$ years ($\approx 1500$ days) gives
standard error $\sqrt{252/1500} \approx 0.41$, so a $95\%$ CI of
$\pm 0.82$. Walk-forward spreads test points over the whole history
and catches regime drift by refitting.

\paragraph{Why it is not enough.} Walk-forward fixes parameter-level
look-ahead, but two leaks remain. (i)~Grid-searching over $100$
values of $W$ and reporting the winner brings Lie~3 back: the
loop was honest, but the choice of $W$ was in-sample. (ii)~It
misses \textbf{preprocessing leakage}: normalising features with
the global mean/variance before the loop leaks information into
every test point. Preprocessing must live inside the loop;
grid-search-on-$W$ is handled by CPCV plus deflated Sharpe, next.

% =====================================================================
\section{Combinatorial Purged Cross-Validation}
\label{sec:bt_cpcv}

Walk-forward is a \emph{single path}: each observation is in the
test set once, in a fixed position. A regime transition landing at
a fit boundary makes the strategy lucky or unlucky, and the reported
Sharpe is one draw from a high-variance distribution.

\textbf{Combinatorial Purged Cross-Validation} (CPCV), due to
\citet{lopezdeprado2018afml} (Ch.~12), generates \emph{many} distinct
train/test splits respecting temporal order and averages the test
statistic. Three principles:
\begin{enumerate}
  \item \textbf{Purging.} Remove from the training set any
  observation whose label-generation window overlaps with test
  observations. This breaks the most insidious leak in ML-based
  trading.
  \item \textbf{Embargo.} Insert a buffer period of $h$
  observations \emph{after} each test fold during which no
  training observations are used. This breaks leakage through
  autocorrelation in the features.
  \item \textbf{Combinatorial splits.} Instead of the usual
  $k$-fold CV's $k$ splits, CPCV uses every combination of two
  (or more) folds as the test set, yielding $\binom{k}{2}$ (or
  more) splits per value of $k$.
\end{enumerate}
First, why naive $k$-fold CV is wrong for time series.

\subsection{Why naive $k$-fold CV is wrong for time series}
\label{subsec:bt_why_kfold_wrong}

Textbook $k$-fold partitions into $k$ folds, trains on $k-1$, tests
on the remainder, repeating $k$ times. Sound for
\emph{i.i.d.}\ data: train and test are independent draws.

Financial time series break i.i.d.\ in two ways:
\begin{enumerate}
  \item \textbf{Feature autocorrelation.} Rolling features (moving
  averages, vols, $z$-scores) are highly correlated between
  neighbours. If $t$ is in test and $t - 1$ in train, the model
  has already seen essentially the test features.
  \item \textbf{Label leakage.} ML-trading labels are often
  forward-looking: ``sign of return over next $H$ bars'' or
  triple-barrier outcomes over $h$ bars. If $t$ is in train and
  $t + 1$ in test, the training label at $t$ was computed from
  the same price path that generates the test features at
  $t + 1$; the model is trained on test data through the
  labels.
\end{enumerate}
The second is fatal: every naive $k$-fold Sharpe in ML-for-trading
is inflated by an amount proportional to the label horizon,
invisible in standard diagnostics. Purging and embargo remove it.

\subsection{Purging and embargo, in words}
\label{subsec:bt_purge_embargo_words}

\textbf{Purging} drops any training observation whose \emph{label
window} overlaps a test observation. If labels at $t$ use price data
in $[t, t+h]$, a training $t$ whose label window intersects a test
range must go: its label was computed from data the model will soon
see as a feature.

\textbf{Embargo} drops training observations for $h$ observations
\emph{after} every test fold. This is feature-side insurance: even
without label overlap, a training point near the test edge may have
highly correlated features.

For test fold $[a, b]$, the surviving training set is
\[
\text{train} \;=\;
  \bigl\{\,t \;:\; [\,t,\, t + h\,] \cap [\,a,\, b\,]
        = \varnothing
        \;\;\text{and}\;\; t \notin [\,b,\, b + h\,]\,\bigr\}.
\]

\subsection{A fully worked 20-period toy example}
\label{subsec:bt_cpcv_toy}

$n = 20$ observations $i = 0, \ldots, 19$; $k = 5$ folds of size $4$;
embargo $h = 1$; combinatorial splits choose $2$ of $5$ folds as the
test set, for $\binom{5}{2} = 10$ splits. Every number below was
verified by the reference implementation at the end of the chapter.

\paragraph{Step 1: folds.} With $n = 20$ and $k = 5$, the fold
boundaries are
\begin{align*}
F_0 &= \{0,1,2,3\}, &
F_1 &= \{4,5,6,7\}, &
F_2 &= \{8,9,10,11\}, \\
F_3 &= \{12,13,14,15\}, &
F_4 &= \{16,17,18,19\}.
\end{align*}

\paragraph{Step 2: the $\binom{5}{2} = 10$ splits.}
The combinatorial splits are all $2$-subsets of $\{0,1,2,3,4\}$:
\begin{align*}
&\underbrace{\{F_0,F_1\}}_{\text{split }1},\;
\{F_0,F_2\},\;
\{F_0,F_3\},\;
\{F_0,F_4\},\;
\{F_1,F_2\},\\
&\{F_1,F_3\},\;
\{F_1,F_4\},\;
\{F_2,F_3\},\;
\{F_2,F_4\},\;
\{F_3,F_4\}.
\end{align*}
So there are $10$ distinct backtest runs instead of the $5$ that a
standard $k$-fold would give you.

\paragraph{Step 3: purge + embargo on split~1, $\{F_0, F_1\}$.}
Test indices: $\{0,1,2,3,4,5,6,7\}$. With an embargo of $h = 1$
after the test set, training observation $i$ must satisfy
\[
i \notin \{0,\ldots,7\}
  \quad \text{and}\quad
  |i - j| > 1 \;\;\text{for every}\;\; j \in \{0,\ldots,7\}.
\]
Of $i = 8, \ldots, 19$, only $i = 8$ is within distance $1$ of
$j = 7$, so it is embargoed. Training indices:
$\{9, 10, \ldots, 19\}$, size $11$. The single $i = 8$ is neither
trained on nor tested on: it is purged.

\paragraph{Step 4: purge + embargo on split~2, $\{F_0, F_2\}$.}
Test indices: $\{0,1,2,3,8,9,10,11\}$, two non-contiguous blocks,
with the embargo applied around each:
\begin{itemize}
  \item Observation $i = 4$ is distance $1$ from $j = 3$
  (embargoed).
  \item Observation $i = 7$ is distance $1$ from $j = 8$
  (embargoed).
  \item Observation $i = 12$ is distance $1$ from $j = 11$
  (embargoed).
\end{itemize}
Training set: $\{5, 6, 13, 14, 15, 16, 17, 18, 19\}$, size $9$.
Three observations ($4, 7, 12$) are purged: the non-contiguous
test set costs more than the contiguous split~1.

\paragraph{Step 5: purge + embargo on split~3, $\{F_0, F_3\}$.}
Test indices: $\{0,1,2,3,12,13,14,15\}$. By the same logic, $i = 4,
11, 16$ are embargoed. Training set:
$\{5, 6, 7, 8, 9, 10, 17, 18, 19\}$, size $9$.

The other seven splits use the same mechanics. Across all $10$, the
researcher gets $10$ out-of-sample P\&L paths. The reported statistic
is the \emph{distribution} of the ten Sharpes, revealing dependence
on which observations fell where. Sharpes
$\{1.3, 1.4, 1.2, 1.5, 1.4, 1.3, 1.2, 1.5, 1.3, 1.4\}$ is stable;
$\{2.8, 1.1, -0.4, 2.2, 0.3, 1.6, 0.8, 2.1, 1.4, -0.2\}$ has the
same mean ($\approx 1.17$) but is one unlucky draw from unrunnable.

\paragraph{Test coverage.} Each of the $5$ folds appears in
$\binom{4}{1} = 4$ of the $10$ splits, so every observation is a
test point $4$ times. Total test-observation count:
$10 \times 8 = 80$, i.e.\ $4$ per observation, four times the
coverage of ordinary $5$-fold, with $10$ (overlap-dependent) Sharpe
estimates.

\begin{keyfactbox}[CPCV --- the three pillars]
Combinatorial Purged CV rests on exactly three rules.
\begin{enumerate}
  \item \textbf{Purge}: drop any training observation whose
  label window overlaps a test observation.
  \item \textbf{Embargo}: drop training observations within
  $h$ of any test observation, buffering feature-side leakage.
  \item \textbf{Combinatorial}: use every $2$-subset of the
  $k$ folds as the test set, yielding $\binom{k}{2}$
  (instead of $k$) train/test splits, each observation
  appearing as a test point $k - 1$ times.
\end{enumerate}
The reported test statistic is the \emph{distribution} across
splits, not any single run's number.
\end{keyfactbox}

\begin{intuitionbox}
CPCV is the minimum cross-validation discipline for time-series ML.
Anything less (naive $k$-fold, single walk-forward, plain train/test)
trains on test data through the labels, with contamination
proportional to the label horizon. Each rule patches a specific
leak: purging $\to$ label leak; embargo $\to$ feature-autocorrelation
leak; combinatorial $\to$ single-path luck. Together they yield a
distribution of out-of-sample Sharpes on disjoint data.
\end{intuitionbox}

\begin{pitfallbox}[Don't skip the label-horizon step]
The common CPCV mistake is implementing combinatorial and embargo
but forgetting that the \emph{label horizon} $h$ is a property of
the label, not the features. If labels at $t$ use prices up to
$t + 20$, then $h = 20$ and any training $t$ whose label window
intersects a test observation within $[t, t+20]$ must be purged.
Purging only $\pm 1$ (confusing embargo for purge) still leaks.
See \citet{lopezdeprado2018afml} Ch.~7.
\end{pitfallbox}

% =====================================================================
\section{Deflated Sharpe Ratio}
\label{sec:bt_deflated_sharpe}

The second statistical lie is multiple-testing inflation. CPCV
controls leakage \emph{inside} a single strategy's evaluation;
deflated Sharpe controls inflation from having tried \emph{many}
strategies. The two are orthogonal; you need both. ``Deflation''
is the plain-English operation of shrinking a reported Sharpe to
account for the number of strategies tried: the best of $100$
zero-edge coin-flip strategies will itself look like an edge, and
the deflation formula removes exactly that bias.

The construction is due to \citet{bailey2014deflated}, summarised in
\citeAFML{} Ch.~14. Under a zero-edge null over $N$ independent
trials, the winner's Sharpe is centred not on zero but on the
\textbf{expected maximum} of $N$ i.i.d.\ Sharpe estimates.

\paragraph{Expected-maximum formula.} For $N$ independent trials
where each Sharpe estimator is asymptotically normal with standard
error $\text{Var}[\widehat{\text{SR}}]^{1/2}$, the expected
maximum is
\begin{equation}
\E\!\left[\max_{n \leq N} \widehat{\text{SR}}_n\right]
  \;\approx\;
  \text{Var}[\widehat{\text{SR}}]^{1/2}
  \Bigl[\,(1 - \gamma)\,\Phi^{-1}\!\bigl(1 - \tfrac{1}{N}\bigr)
       + \gamma\,\Phi^{-1}\!\bigl(1 - \tfrac{1}{N\mathrm{e}}\bigr)\Bigr],
  \label{eq:dsr_emax}
\end{equation}
where $\gamma \approx 0.5772$ is Euler--Mascheroni and $\Phi$ is the
standard normal CDF \citep{bailey2014deflated}. The bracket evaluates
to $\approx 2.5306$ ($N=100$), $3.2414$ ($N=1000$), $3.8513$
($N=10000$). The asymptotic is $\E[\max] \approx \sqrt{2 \ln N}$.

\paragraph{Deflated Sharpe.} The deflated Sharpe ratio $\text{DSR}$
is then defined as the probability, under the null, that the
observed Sharpe $\widehat{\text{SR}}$ exceeds the expected maximum
under the best of $N$ trials:
\begin{equation}
\text{DSR}
  \;\defeq\;
  \Phi\!\left(
    \frac{\bigl(\widehat{\text{SR}} - \E[\max \widehat{\text{SR}}]\bigr)
          \sqrt{T - 1}}
         {\sqrt{1 - \hat\gamma_3\,\widehat{\text{SR}}
                + \tfrac{\hat\gamma_4 - 1}{4}\widehat{\text{SR}}^2}}
    \right),
  \label{eq:dsr_definition}
\end{equation}
where $T$ is the number of return observations in the Sharpe
estimate, and $\hat\gamma_3$, $\hat\gamma_4$ are the sample skewness
and kurtosis of the return series \citep{bailey2014deflated}. Under
normal returns ($\hat\gamma_3 = 0$, $\hat\gamma_4 = 3$) the
denominator simplifies to $\sqrt{1 + \widehat{\text{SR}}^2/2}
\approx 1$ for small $\widehat{\text{SR}}$.

Convention: $\text{DSR} \geq 0.95$ means the strategy is unlikely to
be the winner of a random $N$-trial tournament; $\text{DSR} \approx
0.50$ is noise.

\subsection{A worked DSR example}
\label{subsec:bt_dsr_worked}

$N = 100$ configurations on five years of daily data, so
$T = 252 \times 5 = 1260$; observed annualised Sharpe of the winner
$\widehat{\text{SR}}_{\text{ann}} = 1.8$; assume normal returns. In
Python, to four decimals:

\begin{enumerate}
  \item Convert to per-period Sharpe:
  $\widehat{\text{SR}} = 1.8 / \sqrt{252} = 0.1134$.
  \item Expected-maximum bracket in
  \eqref{eq:dsr_emax} for $N = 100$:
  $(1 - 0.5772)\Phi^{-1}(1 - 0.01)
   + 0.5772\,\Phi^{-1}(1 - 1/(100\mathrm{e}))
   = 0.4228 \times 2.3263 + 0.5772 \times 2.5894
   = 0.9836 + 1.4946
   = 2.4782.$
  (A more careful computation yields $2.5306$; the small
  discrepancy is from higher-order terms in the approximation.)
  \item Multiply by $\text{Var}[\widehat{\text{SR}}]^{1/2}
  = 1/\sqrt{T} = 1/\sqrt{1260} = 0.02817$ to get the expected
  maximum per-period Sharpe:
  $\E[\max \widehat{\text{SR}}] \approx 2.5306 \times 0.02817
  = 0.07129$.
  \item Plug into the denominator of
  \eqref{eq:dsr_definition}: under normal returns
  ($\hat\gamma_3 = 0$, $\hat\gamma_4 = 3$) the denominator is
  $\sqrt{1 + 0.1134^2 / 2}/\sqrt{T - 1}
   = \sqrt{1.00643}/\sqrt{1259}
   = 0.02827$.
  \item Compute the DSR argument:
  $(0.1134 - 0.07129)/0.02827 = 1.4889$.
  \item Evaluate $\Phi$: $\Phi(1.4889) = 0.9317$.
\end{enumerate}
DSR $\approx 0.93$: under $N = 100$, a $93\%$ chance of genuine edge
rather than tournament noise. At $N = 1000$ the expected maximum
climbs to $3.2414 \times 0.02817 = 0.0913$, argument drops to
$(0.1134 - 0.0913)/0.02827 = 0.7831$, $\text{DSR} = \Phi(0.7831) =
0.7832$, below $0.95$. Ten times more trials flips ``probably
real'' to ``probably noise'' with no data change.

\paragraph{Headline.} The \emph{effective} annualised Sharpe is
\[
\widehat{\text{SR}}_{\text{eff}}
  \;=\; \widehat{\text{SR}} - \E[\max \widehat{\text{SR}}],
\]
i.e.\ $1.8 - 1.13 = 0.67$ for $N=100$. A reported Sharpe of $1.8$
from $100$ trials is worth $\approx 0.67$ after deflation.

\begin{keyfactbox}[Deflated Sharpe Ratio]
Under $N$ independent trials with per-period Sharpe standard
error $1/\sqrt{T}$, the expected maximum Sharpe under a
zero-edge null is approximately $\E[\max] / \sqrt{T}$ where
$\E[\max]$ is the expected maximum of $N$ iid standard normals.
The effective (deflated) annualised Sharpe is
\[
\widehat{\text{SR}}_{\text{eff}}
  \;\approx\;
  \widehat{\text{SR}}
  - \E[\max] \cdot \sqrt{\frac{252}{T}}.
\]
For $\widehat{\text{SR}} = 1.8$, $T = 1260$, $N = 100$ (using
the Bailey bracket $\E[\max] \approx 2.5306$):
$\widehat{\text{SR}}_{\text{eff}} \approx 1.8 - 1.13 = 0.67$.
Ten times more trials ($N = 1000$, $\E[\max] \approx 3.2414$)
reduces $\widehat{\text{SR}}_{\text{eff}}$ to
$\approx 1.8 - 1.45 = 0.35$.
\end{keyfactbox}

\begin{pitfallbox}[Multiple testing silently inflates Sharpe]
Six months of research typically touches $N \in [50, 500]$ trials.
Report the winner's Sharpe without tracking $N$ and the number is
inflated by $\sqrt{2 \ln N}$ standard errors; the true effective
Sharpe is $\widehat{\text{SR}} - \sqrt{2 \ln N \cdot 252 / T}$. For
$T = 1260$ and that range, subtract $0.88$--$1.41$. Any
$\widehat{\text{SR}} \leq 1.4$ from six months of research is
indistinguishable from chance at $N = 500$; do not deploy without
(a)~extending the sample, (b)~a fresh out-of-sample hold-out, or
(c)~live tiny-capital paper trading.
\end{pitfallbox}

% =====================================================================
\section{Backtesting-framework comparison}
\label{sec:bt_frameworks}

Four Python frameworks dominate.

\paragraph{\texttt{Backtesting.py}} (Zvezdin et al.). Compact,
event-driven, single-asset. Small API, clean docs, $20$-line
strategies, simple cost/slippage, built-in hyperparameter
optimiser. No built-in walk-forward or CPCV, no portfolio
constraints, no order book. Good for learning.

\paragraph{\texttt{vectorbt}} (Polakow). Pandas-native, fully
vectorised, optimised for parameter sweeps. A $10{,}000$-point
MA-crossover grid runs in seconds. Supports portfolios, sizing,
slippage. Not truly event-driven: signals and fills align bar by
bar, so intra-bar order-book dynamics are out of scope, and the
vector API tempts next-bar look-ahead.

\paragraph{\texttt{NautilusTrader}} (Nautech Systems). Rust core,
Python API, tick-by-tick with a full order-book model: it can
replay raw LOBSTER dumps with realistic queue-position fills.
Operationally heavy; overkill for daily-bar sweeps. Use when
microstructure correctness matters.

\paragraph{\texttt{QuantConnect} / \texttt{Lean}.} Cloud-hosted
full-stack platform with point-in-time historical data for US
equities, options, futures, crypto, and an event-driven engine.
The PIT universes are the main reason to pay. Lock-in: strategies
are written to the Lean API. Prosperity does not need it; GS desks
use internal equivalents.

\begin{table}[ht]
\centering
\footnotesize
\begin{tabular}{@{}lllll@{}}
\toprule
Framework & Model & Speed & CV support & Best for \\
\midrule
\texttt{Backtesting.py}
  & event, 1 asset
  & moderate
  & manual wrap
  & learning \\
\texttt{vectorbt}
  & bar-aligned vector
  & very fast
  & manual wrap
  & parameter sweeps \\
\texttt{NautilusTrader}
  & event, full book
  & slow
  & manual wrap
  & microstructure \\
\texttt{QuantConnect}
  & event, cloud
  & moderate
  & built-in
  & serious PIT data \\
\bottomrule
\end{tabular}
\caption{One-line trade-offs for the four mainstream Python
backtesting frameworks. None ship CPCV out of the box except
QuantConnect, and even there it is a research tool rather than a
core primitive; walk-forward and CPCV should be wrapped around
whichever framework is chosen.}
\label{tab:bt_framework_comparison}
\end{table}

\begin{keyfactbox}[Framework taxonomy]
\begin{itemize}
  \item \texttt{Backtesting.py}: simple, event-driven,
  single-asset. Teach yourself on this.
  \item \texttt{vectorbt}: vectorised, very fast. Best for
  parameter sweeps. Easy to leak look-ahead if you're not
  careful.
  \item \texttt{NautilusTrader}: tick-level, order-book,
  production-grade. Use when microstructure matters.
  \item \texttt{QuantConnect} / \texttt{Lean}: point-in-time
  data, cloud-hosted, institutional-ish. Use when survivorship
  bias is the binding constraint.
\end{itemize}
\end{keyfactbox}

% =====================================================================
\section{Prosperity case study: plateau, not peak}
\label{sec:bt_prosperity}

\begin{prosperitybox}
The Frankfurt Hedgehogs, 2nd in Prosperity~3, describe their
parameter-selection discipline (see
\texttt{imc-prosperity-3/README.md}, \S$410$):

\begin{quote}
``During parameter selection, we always prioritized landscape
stability over pure performance peaks. Rather than picking the
best parameter set based on maximum backtested profit, we
chose combinations that showed consistent, flat regions of
good performance, reducing sensitivity to slight shifts and
avoiding overfit disasters.''
\end{quote}

In this chapter's language: they reported the top of a \emph{plateau}
(a neighbourhood where Sharpe is roughly constant), not an in-sample
peak. Plateaus defend against Lies~3, 4, and 7 at once: a narrow
spike usually fits noise (a $10\%$ perturbation of the lookback kills
the edge), a flat region is edge with a margin of safety.

Implementation: instead of \texttt{argmax} over the grid, return the
point whose $\pm k$-neighbourhood average is maximised. For two
parameters $(W, z_{\text{entry}})$ a $3 \times 3$ window suffices.
\end{prosperitybox}

% =====================================================================
\section{Production pipeline on a sell-side desk}
\label{sec:bt_gs}

\begin{gsbox}
On a GS systematic desk, a new strategy clears five gates in order:

\begin{enumerate}
  \item \textbf{Walk-forward on $\geq 5$ years.} The Strat owns
  the loop. Five years is the minimum to span two distinct
  macro regimes at daily frequency.
  \item \textbf{Deflated Sharpe.} Report raw Sharpe, trial
  count, and DSR per \eqref{eq:dsr_definition}. DSR below
  $0.95$ kills the strategy.
  \item \textbf{Transaction-cost audit.} The execution desk
  reprices every trade using the in-house impact model
  (calibrated square-root), adds financing/clearing, and
  returns a net Sharpe. The most common casualty.
  \item \textbf{Paper trading $\geq 30$ days} on a production
  router with orders intercepted before the exchange. Paper
  vs.\ projected P\&L divergence above $\sim 1$ Sharpe point
  or $20\%$ of return triggers review.
  \item \textbf{Risk limits, margin, kill-switches.} The risk
  system (\cref{ch:risk_production}) owns this gate: declared
  max notional, gross, net, Greeks, and auto-liquidation
  trigger.
\end{enumerate}
A Strat who skips gate~3 and reports Sharpe $2.0$ to the desk head
will be asked: ``Net of execution costs, what is the Sharpe?''
Always have the number ready.
\end{gsbox}

% =====================================================================
\section{A short reference implementation}
\label{sec:bt_reference_impl}

A minimal CPCV splitter fits in under $30$ lines. The implementation
below matches \cref{subsec:bt_cpcv_toy}: \texttt{n=20, k=5,
embargo=1} reproduces the $10$ splits, including purged $i = 8$ on
split~$1$. Uses $h = 1$ for clarity; production extends the inner
check to arbitrary label windows.

\begin{lstlisting}[language=Python, caption={Minimal CPCV splitter,
verified against the toy example in this chapter.}]
import numpy as np
from itertools import combinations

def cpcv_splits(n, k, embargo):
    """Generate (train_idx, test_idx) pairs under combinatorial purged CV.

    n:       number of observations
    k:       number of folds (test sets will be all 2-subsets of folds)
    embargo: number of observations to drop on each side of the test set
    """
    fold_size = n // k
    all_folds = [list(range(i * fold_size, (i + 1) * fold_size))
                 for i in range(k)]
    for test_fold_indices in combinations(range(k), 2):
        test_idx = sum([all_folds[i] for i in test_fold_indices], [])
        test_set = set(test_idx)
        train_idx = [i for i in range(n)
                     if i not in test_set
                     and not any(abs(i - j) <= embargo for j in test_idx)]
        yield train_idx, test_idx

# Sanity-check against the 20/5/1 toy example in the chapter:
splits = list(cpcv_splits(n=20, k=5, embargo=1))
assert len(splits) == 10
train0, test0 = splits[0]  # test folds (F_0, F_1)
assert test0 == [0, 1, 2, 3, 4, 5, 6, 7]
assert train0 == [9, 10, 11, 12, 13, 14, 15, 16, 17, 18, 19]  # 8 is purged
\end{lstlisting}

The two asserts are the minimum regression test: split count
$\binom{5}{2} = 10$ and purged $i = 8$ on the first split. Passing
both does not guarantee correctness for arbitrary label horizons,
but confirms the purge-plus-embargo logic on the simple case. For
$h > 1$, replace \texttt{any(abs(i - j) <= embargo ...)} with a
label-window intersection test per \citeAFML{} Ch.~7.

% =====================================================================
\section{Key facts to memorise}
\label{sec:bt_key_facts}

\begin{itemize}
  \item \textbf{The seven lies}: look-ahead, survivorship,
  overfitting, data snooping, transaction-cost neglect,
  slippage, regime dependence. Every backtest issue is one of
  these or a composition.
  \item \textbf{Walk-forward}: fit on $[t-W, t)$, trade on
  $[t, t+R)$, concatenate test P\&Ls. Fixes most of Lie~1;
  not Lies~3 or 4.
  \item \textbf{Naive $k$-fold is wrong for time series}:
  labels over horizon $h$ leak into training even when
  features do not.
  \item \textbf{CPCV pillars}: purge (drop training whose
  label window intersects test); embargo (drop training
  within $h$ of the test fold); combinatorial (all
  $\binom{k}{2}$ two-fold splits). Report a
  \emph{distribution}.
  \item \textbf{CPCV coverage}: each observation is a test
  point in $k - 1$ of the $\binom{k}{2}$ splits, $k - 1$
  times richer than $k$-fold.
  \item \textbf{Expected max Sharpe (null)} per-period:
  $\E[\max]/\sqrt{T}$, $\E[\max]$ the max of $N$ iid
  normals. Bailey: $2.53$ ($N=100$), $3.24$ ($N=1000$),
  $3.85$ ($N=10000$).
  \item \textbf{DSR headline}: $\widehat{\text{SR}} = 1.8$,
  $T = 1260$, $N = 100$ gives $\widehat{\text{SR}}_{\text{eff}}
  \approx 0.67$; $N = 1000$ gives $\approx 0.35$.
  \item \textbf{Costs are life-or-death}: gross Sharpe $1.59$
  on $5$~bp daily edge becomes $0.95$ at $2$~bp round-trip
  and $0$ at $5$~bp.
  \item \textbf{Framework taxonomy}: \texttt{Backtesting.py}
  (learning), \texttt{vectorbt} (speed),
  \texttt{NautilusTrader} (microstructure),
  \texttt{QuantConnect} (PIT data). None ship CPCV; wrap.
  \item \textbf{Plateau not peak}: prefer a flat-region
  parameter to a spike. Frankfurt Hedgehogs' Prosperity~3
  lesson.
  \item \textbf{GS gates}: walk-forward $\geq 5$ years,
  DSR, execution cost audit, $30$-day paper, risk approval.
  \item \emph{A backtest is a hypothesis; walk-forward, CPCV,
  and deflated Sharpe are what make it falsifiable.} Without
  them it is a story.
\end{itemize}
